\chapter{Assets, Authority, and Trust Boundaries}

\section{Purpose}

Chapter 1 framed blockchain systems as adversarial state machines.
Chapter 2 asks the next question:

\begin{quote}
Which parts of the state are valuable, and who has the authority to
change them?
\end{quote}

Most blockchain security failures involve confusion about assets,
authority, or trust boundaries. A protocol team says ``users own their
funds,'' but an upgrade key can replace the withdrawal logic. A token
claims to be backed by off-chain reserves, but the reserve proof is only
an attestation. A bridge mints a wrapped asset, but the wrapped asset is
only as strong as the source-chain finality, validator set, relayers,
and message verification path. A governance system looks decentralized,
but a small quorum, delegated voting, emergency guardian, or timelock
admin can still decide outcomes.

Security analysis becomes unserious when it treats ``ownership'' as a
word rather than a state-machine fact.

For this book, ownership means nothing until you can answer:

\begin{lstlisting}
Which state variable, object, output, role, or claim records the asset?
Which transition moves or modifies it?
Which authority can cause that transition?
Which assumptions make that authority trustworthy enough?
\end{lstlisting}

This chapter builds the vocabulary for those answers.

\section{Learning Objectives}

By the end of this chapter, you should be able to:

\begin{itemize}
\tightlist
\item
  Define on-chain assets as enforceable claims over state, not just
  tokens.
\item
  Distinguish ownership, possession, authority, custody, control, and
  economic exposure.
\item
  Compare ownership models across UTXOs, accounts, contracts, Solana
  accounts, and Move resources.
\item
  Identify privileged authorities such as admins, governors, validators,
  sequencers, custodians, oracles, relayers, and upgrade controllers.
\item
  Draw trust boundaries across on-chain code, off-chain infrastructure,
  human processes, legal claims, and cross-chain systems.
\item
  Build an authority map showing who can move, freeze, mint, burn,
  upgrade, censor, or otherwise affect assets.
\end{itemize}

\section{The Core Distinction}

An asset is something valuable represented by, controlled by, or
dependent on system state.

Authority is the ability to cause state transitions that affect that
asset.

These are related, but not the same.

An address may hold tokens but have granted a spender allowance. A user
may own vault shares but depend on a strategist, oracle, withdrawal
queue, and admin pause key. A DAO may own a treasury but a timelock
executor may be able to trigger the transfer. A stablecoin holder may
own an on-chain balance but depend on off-chain reserve custody and
redemption policy. A rollup user may own an L2 balance but depend on the
L1 bridge contract, proof system, data availability, and exit path.

The basic security questions are:

\begin{lstlisting}
Asset:
  What has value?

State:
  Where is that value represented?

Authority:
  Who or what can change that representation?

Boundary:
  Which facts must be trusted rather than verified locally?
\end{lstlisting}

If you cannot answer those questions, you do not understand the system's
security model yet.

\section{What Counts as an Asset On-Chain}

\subsection{Assets Are Enforced Claims Over State}

An on-chain asset is not a tiny object sitting inside an address. It is
a claim recognized by a state-transition system.

For example:

\begin{itemize}
\tightlist
\item
  A Bitcoin coin is spendable because an unspent transaction output
  exists and a future transaction can satisfy its spending condition.
\item
  ETH is spendable because an Ethereum account balance exists and a
  valid transaction can debit it.
\item
  An ERC-20 token balance is spendable because a contract's state maps
  an address to a number and the contract's transfer rules allow that
  number to decrease and another to increase.
\item
  An ERC-721 NFT is transferable because a contract maps a token ID to
  an owner and its transfer rules recognize the caller as owner or
  authorized operator.
\item
  A Solana token account records token state inside an account whose
  data is governed by a program.
\item
  A Move resource represents value partly because the language and
  module system restrict copying, dropping, storing, and unpacking.
\item
  A bridged asset is a claim in one domain that depends on locked,
  burned, proven, or otherwise accounted value in another domain.
\item
  A governance token is both an economic asset and a control surface.
\end{itemize}

The asset exists because the system's rules say certain future
transitions are allowed and others are not.

That is why security review cannot stop at ``who owns the token.'' The
real question is:

\begin{lstlisting}
Which transitions can affect the claim?
\end{lstlisting}

\subsection{Asset Categories}

A useful asset inventory should include more than obvious balances.

Common blockchain asset categories:

{\def\LTcaptype{none} % do not increment counter
\begin{longtable}[]{@{}
  >{\raggedright\arraybackslash}p{(\linewidth - 4\tabcolsep) * \real{0.3333}}
  >{\raggedright\arraybackslash}p{(\linewidth - 4\tabcolsep) * \real{0.3333}}
  >{\raggedright\arraybackslash}p{(\linewidth - 4\tabcolsep) * \real{0.3333}}@{}}
\toprule\noalign{}
\begin{minipage}[b]{\linewidth}\raggedright
Category
\end{minipage} & \begin{minipage}[b]{\linewidth}\raggedright
Examples
\end{minipage} & \begin{minipage}[b]{\linewidth}\raggedright
Security question
\end{minipage} \\
\midrule\noalign{}
\endhead
\bottomrule\noalign{}
\endlastfoot
Native assets & BTC, ETH, SOL, staking rewards & Who can spend, slash,
censor, or reorder transactions affecting them? \\
Fungible tokens & ERC-20, SPL tokens, stablecoins & Who can transfer,
approve, mint, burn, freeze, or blacklist? \\
Non-fungible assets & ERC-721 tokens, game items, account-bound records
& Who controls transfer, metadata, approvals, and operator authority? \\
Semi-fungible assets & ERC-1155-style assets, batch-minted items & Are
balances, IDs, and batch operations checked consistently? \\
Shares and claims & LP tokens, vault shares, receipt tokens, restaking
positions & What do shares claim, and can the exchange rate be
manipulated? \\
Debt and liabilities & Borrow positions, margin accounts,
undercollateralized debt & Who can create, liquidate, forgive, or
socialize debt? \\
Capabilities and roles & Admin roles, minter roles, pause guardians,
upgrade keys & What privileged transitions do they unlock? \\
Governance power & Voting tokens, delegation, proposal rights & Can
governance change the rules or seize assets? \\
Cross-domain claims & Wrapped assets, bridge receipts, rollup
withdrawals & Which external state and finality assumptions are
trusted? \\
Off-chain claims & Fiat-backed stablecoins, RWAs, redemption rights &
Which legal, custody, and reporting systems are outside the chain? \\
Information assets & Oracle prices, validator keys, private order flow,
signer policies & Who can distort or leak information that controls
value? \\
\end{longtable}
}

The point of this table is not classification for its own sake. It is to
force a reviewer to notice assets that are not called ``tokens.''

A governance role is an asset. An oracle key is an asset. A bridge
validator key is an asset. An unexecuted withdrawal claim is an asset. A
pending governance proposal may become an asset or a liability depending
on who can execute it.

\subsection{Native Assets and Contract Assets Differ}

Native assets are accounted for by the base protocol. Contract assets
are accounted for by contract or program logic.

On Ethereum, ETH balances are part of account state, while ERC-20
balances are state inside a token contract. Ethereum's documentation
describes accounts as entities with ETH balances that can send messages,
and distinguishes externally owned accounts controlled by private keys
from contract accounts controlled by code.\footnote{ethereum.org,
  ``Ethereum Accounts'', https://ethereum.org/developers/docs/accounts/.}

That distinction matters.

If a user sends ETH, the base protocol updates account balances
according to transaction rules. If a user transfers an ERC-20 token, the
token contract updates its internal balance mapping according to its own
implementation. The ERC-20 standard specifies functions such as
\passthrough{\lstinline!balanceOf!}, \passthrough{\lstinline!transfer!},
\passthrough{\lstinline!transferFrom!},
\passthrough{\lstinline!approve!}, and
\passthrough{\lstinline!allowance!}, but different implementations can
still behave differently within or around that interface.\footnote{Fabian
  Vogelsteller and Vitalik Buterin, ``ERC-20: Token Standard'',
  https://eips.ethereum.org/EIPS/eip-20.}

Security consequences:

\begin{itemize}
\tightlist
\item
  A wallet may display both ETH and ERC-20 balances, but those balances
  are enforced by different rule systems.
\item
  A transfer of ETH and a transfer of a token can have different failure
  behavior.
\item
  A token can have minting, burning, pausing, blacklisting, rebasing,
  transfer fees, hooks, or nonstandard return values.
\item
  A protocol that treats all tokens as if they were identical ERC-20s
  will eventually be wrong.
\end{itemize}

Do not say ``the user has funds'' until you know which accounting system
recognizes the claim.

\subsection{An Asset Can Be a Right, Not a Balance}

Many important assets are rights to do something later.

Examples:

\begin{itemize}
\tightlist
\item
  A withdrawal ticket.
\item
  A bridge message that has been proven but not executed.
\item
  A governance proposal that can be queued.
\item
  A validator withdrawal credential.
\item
  A liquidation right.
\item
  A claim on future rewards.
\item
  A timelocked allocation.
\item
  A vesting stream.
\item
  A pending unstake.
\item
  A signature that authorizes a permit or order.
\end{itemize}

These rights may not look like balances, but they are still valuable.
Their security depends on whether they can be created, transferred,
consumed, canceled, replayed, or blocked.

For every such right, ask:

\begin{lstlisting}
How is the right represented?
Who can create it?
Who can consume it?
Can it be consumed more than once?
Can it expire?
Can it be canceled?
Can it be transferred?
Can it be made unusable?
\end{lstlisting}

Rights are especially important in cross-chain, lending, governance, and
staking systems, where value often exists in an intermediate state.

\subsection{Off-Chain Claims Are Not Made On-Chain by Naming Them}

Some assets are on-chain records of off-chain claims.

Examples:

\begin{itemize}
\tightlist
\item
  Fiat-backed stablecoins.
\item
  Tokenized treasury bills.
\item
  Real-world asset tokens.
\item
  Custodial wrapped BTC.
\item
  Claims on carbon credits, invoices, real estate, or commodities.
\end{itemize}

The on-chain token may be transferred by code, but the underlying claim
depends on legal agreements, custody, redemption policies, banking
access, auditors, administrators, and sometimes court systems.

The chain can enforce:

\begin{lstlisting}
address A has 100 tokens
address A can transfer 100 tokens if not blocked
\end{lstlisting}

The chain usually cannot enforce:

\begin{lstlisting}
the issuer has matching reserves
the custodian did not rehypothecate assets
the bank account remains accessible
the token holder has legally enforceable redemption rights
the reserve report is complete and current
\end{lstlisting}

This does not make off-chain claims worthless. It means they cross a
trust boundary. A serious security analysis states that boundary
directly.

Bad:

\begin{lstlisting}
The token is backed 1:1.
\end{lstlisting}

Better:

\begin{lstlisting}
The token contract records transferable claims. Redemption depends on issuer solvency, reserve custody, banking access, legal enforceability, and the issuer's honoring of redemptions.
\end{lstlisting}

The second statement is less marketable. It is also more honest.

\subsection{Control Surfaces Are Assets}

Attackers often target control before they target balances.

Control surfaces include:

\begin{itemize}
\tightlist
\item
  Upgrade authority.
\item
  Mint authority.
\item
  Burn authority.
\item
  Pause authority.
\item
  Blacklist or freeze authority.
\item
  Oracle update authority.
\item
  Governance proposal rights.
\item
  Timelock proposer, executor, and canceller roles.
\item
  Bridge validator or relayer keys.
\item
  Sequencer control.
\item
  Validator keys.
\item
  Backend signer keys.
\item
  Deployment keys and CI/CD secrets.
\item
  Frontend domains and package-publishing accounts.
\end{itemize}

These are assets because controlling them can lead to asset movement,
denial of service, governance capture, or reputation damage.

If a protocol has \$10 million in a vault and a single key can upgrade
the vault to drain funds, the key is economically equivalent to control
over the vault. Treat it accordingly.

\section{Ownership: Keys, Accounts, UTXOs, Objects, and Claims}

\subsection{Ownership Is a State-Machine Relationship}

In ordinary language, ownership sounds legal or moral. In blockchain
security, ownership is operational:

\begin{lstlisting}
An actor owns an asset to the extent that the system recognizes that actor,
or an authority controlled by that actor,
as able to exercise the relevant transitions over the asset.
\end{lstlisting}

This definition is intentionally dry. It prevents category errors.

A private key may control an account. The account may hold a token. The
token may represent a claim on reserves. The reserves may be held by a
custodian. The custodian may be subject to law. The user may have legal
rights against the issuer. These are different layers of ownership and
authority.

Do not collapse them.

\subsection{UTXO Ownership: Spend Conditions, Not Account Balances}

In a UTXO system, the state contains unspent outputs. Each output has a
value and a spending condition. To ``own'' the output is to be able to
satisfy that condition in a valid transaction.

Bitcoin's whitepaper describes coins as transferred through digital
signatures and then immediately identifies the hard problem: the
recipient needs assurance that an earlier owner did not
double-spend.\footnote{Satoshi Nakamoto, ``Bitcoin: A Peer-to-Peer
  Electronic Cash System'', https://bitcoin.org/bitcoin.pdf.} That is
why ownership in Bitcoin is not just ``this public key appears in a
chain of signatures.'' It also depends on the network agreeing which
spend came first.

Security consequences:

\begin{itemize}
\tightlist
\item
  The same key may control many outputs.
\item
  A wallet balance is a derived view over a set of spendable outputs.
\item
  Spending one output usually consumes it and creates new outputs.
\item
  Coin selection affects privacy, fees, and sometimes security.
\item
  Script conditions can express multisig, timelocks, hashlocks, and
  other constraints.
\item
  A UTXO can be unspendable if its script is impossible to satisfy or if
  required keys are lost.
\end{itemize}

The security question is not:

\begin{lstlisting}
Which account owns the money?
\end{lstlisting}

It is:

\begin{lstlisting}
Which unspent outputs exist, and what conditions allow each to be spent?
\end{lstlisting}

\subsection{Account Ownership: Keys, Code, and Nonces}

Account-based systems present a different interface. On Ethereum,
externally owned accounts are controlled by private keys; contract
accounts are controlled by code.\footnote{ethereum.org, ``Ethereum
  Accounts'', https://ethereum.org/developers/docs/accounts/.}

This distinction is foundational.

An EOA can initiate transactions because someone can produce signatures
for it. A contract account cannot initiate transactions on its own in
the same way; it executes when called, and its authority is defined by
code and state.

Security consequences:

\begin{itemize}
\tightlist
\item
  Losing an EOA private key usually means losing practical control over
  the account's assets.
\item
  A compromised EOA can sign valid transactions that are disastrous.
\item
  A contract can hold assets but expose only specific methods for moving
  them.
\item
  A contract can be the owner of another contract.
\item
  Nonces prevent simple replay of transactions from the same account.
\item
  Smart accounts and account abstraction blur the interface, but not the
  underlying question: which validation logic authorizes the transition?
\end{itemize}

Ethereum.org also notes that an account is not a wallet; a wallet is an
interface for interacting with an account.\footnote{ethereum.org,
  ``Ethereum Accounts'', https://ethereum.org/developers/docs/accounts/.}
That distinction becomes important in user-side security. A wallet can
mislead a user, fail to simulate consequences, connect to a malicious
site, or sign the wrong message. The account is the state object; the
wallet is a tool that controls or mediates authority.

\subsection{Contract Ownership Is Just State}

When a Solidity contract has an \passthrough{\lstinline!owner!}, that
owner is usually a state variable checked by privileged functions.

Example:

\begin{lstlisting}
address public owner;

modifier onlyOwner() {
    require(msg.sender == owner, "not owner");
    _;
}

function setFee(uint256 newFee) external onlyOwner {
    fee = newFee;
}
\end{lstlisting}

This is not mystical. It means:

\begin{lstlisting}
State:
  owner
  fee

Transition:
  setFee(newFee)

Validity:
  caller == owner

Effect:
  fee = newFee
\end{lstlisting}

OpenZeppelin's access-control documentation emphasizes that access
control governs who can perform actions such as minting, voting,
freezing transfers, or administrative tasks.\footnote{OpenZeppelin,
  ``Access Control'',
  https://docs.openzeppelin.com/contracts/5.x/access-control.} It also
distinguishes simple ownership from role-based access control, where
separate roles can authorize different actions.

The security reviewer should ask:

\begin{itemize}
\tightlist
\item
  Who is the owner?
\item
  Can ownership be transferred?
\item
  Is transfer one-step or two-step?
\item
  Can ownership be renounced?
\item
  If ownership is renounced, does the system lose necessary recovery
  powers?
\item
  If ownership is transferred to a contract, what controls that
  contract?
\item
  If ownership is transferred to a multisig, who controls the signer
  keys?
\item
  If ownership is transferred to governance, can governance be captured?
\item
  If ownership is transferred to a timelock, who can propose, cancel,
  and execute?
\end{itemize}

An \passthrough{\lstinline!owner!} variable is rarely the end of the
authority chain. It is usually the beginning.

\subsection{Token Ownership Is Defined by the Token's Rules}

Token standards define interfaces, not complete security models.

ERC-20 defines basic functions for balances, transfers, approvals, and
allowances. It explicitly includes delegated spending through
\passthrough{\lstinline!approve!},
\passthrough{\lstinline!transferFrom!}, and
\passthrough{\lstinline!allowance!}.\footnote{Fabian Vogelsteller and
  Vitalik Buterin, ``ERC-20: Token Standard'',
  https://eips.ethereum.org/EIPS/eip-20.} That means a user's token
balance is not the only state relevant to movement. Allowances are
authority too.

For an ERC-20-like token, asset movement may be authorized by:

\begin{itemize}
\tightlist
\item
  The holder calling \passthrough{\lstinline!transfer!}.
\item
  A spender calling \passthrough{\lstinline!transferFrom!} using
  allowance.
\item
  A minter increasing supply.
\item
  A burner reducing supply or destroying a user's balance.
\item
  An admin pausing transfers.
\item
  A blacklist authority blocking addresses.
\item
  A rebasing mechanism changing balances.
\item
  A bridge contract minting or burning representations.
\end{itemize}

ERC-721 adds different authority semantics. The standard tracks
ownership by token ID, includes approvals, and includes operators that
can manage all NFTs of an owner.\footnote{William Entriken et al.,
  ``ERC-721: Non-Fungible Token Standard'',
  https://eips.ethereum.org/EIPS/eip-721.} A security review of NFTs
therefore must inspect not only
\passthrough{\lstinline!ownerOf(tokenId)!} but also per-token approvals,
operator approvals, transfer hooks, metadata authority, and marketplace
integration.

The general rule:

\begin{lstlisting}
Do not infer token authority from the word "owner."
Read the token's transition rules.
\end{lstlisting}

\subsection{Solana Ownership: Account Owner Program and Signers}

Solana uses ``owner'' in a way that can confuse EVM-trained reviewers.

Solana documentation describes accounts as the fundamental data units
storing state, with each account identified by a 32-byte address. It
also states that only the account's owner program can modify its data or
debit lamports, while any program can credit lamports to a writable
account.\footnote{Solana Foundation, ``Accounts'',
  https://solana.com/docs/core/accounts.}

This is not the same as ``the user owns this account'' in ordinary
language.

For this chapter, keep only the ownership lesson: separate the program
that can modify account data, the signer that authorized the
instruction, and the account constraints that prove the passed accounts
are the intended accounts. Part III will study the detailed Solana
execution model.

For Solana, ask:

\begin{lstlisting}
Which program owns this account's data?
Which signer authorized this instruction?
Which accounts are writable?
Which account constraints prove that these are the intended accounts?
Which PDA seeds establish program authority?
\end{lstlisting}

\subsection{Move Ownership: Resources and Module Boundaries}

Move was designed around resource-like values whose abilities control
whether they can be copied, dropped, stored, or used as storage keys.
The current Move Book describes abilities as the typing feature that
controls what actions are permissible for values of a given type,
including linear behavior and storage use.\footnote{The Move Book,
  ``Abilities'', https://move-book.com/reference/abilities/.}

Move changes the shape of some asset bugs. It can make accidental
duplication or dropping of resources harder than in less restrictive
models. It does not remove the need to map creation, destruction,
transfer, privileged capabilities, and module boundaries. Part III will
handle the execution-model details.

You still need to ask:

\begin{itemize}
\tightlist
\item
  Which module defines the resource?
\item
  Which functions can create it?
\item
  Which functions can destroy it?
\item
  Which functions can transfer or store it?
\item
  Which capabilities authorize privileged actions?
\item
  Which invariants are enforced by the type system, and which remain
  protocol-level assumptions?
\end{itemize}

Move's resource model is a strong tool for representing assets. It is
not a substitute for designing correct authority.

\subsection{Legal Ownership and Protocol Control Can Diverge}

Sometimes the person legally entitled to an asset is not the actor who
can move it on-chain.

Examples:

\begin{itemize}
\tightlist
\item
  A user deposits into a centralized exchange. The exchange controls the
  withdrawal wallet.
\item
  A stablecoin holder may have an on-chain balance but redemption
  depends on issuer terms.
\item
  A custodian signs transactions for an institution.
\item
  A DAO treasury may be legally associated with a foundation while
  on-chain control sits with a multisig or governance contract.
\item
  A court order may require freezing assets, but the chain only enforces
  that if some authority can freeze them.
\item
  A compromised key can move assets even though the attacker is not the
  rightful owner.
\end{itemize}

For security work, distinguish:

\begin{lstlisting}
Legal title:
  Who has an off-chain claim?

Protocol authority:
  Who can cause the state transition?

Custody:
  Who controls the signing material or operational process?

Economic exposure:
  Who gains or loses value if the state changes?
\end{lstlisting}

These categories often overlap. When they do not, risk appears.

\section{Authority: Users, Admins, Validators, Governance, and Custodians}

\subsection{Authority Is the Ability to Cause a Transition}

Authority is not status. It is not reputation. It is not what the docs
say someone should do. Authority is the ability to cause a state
transition that the system will accept.

An authority may be:

\begin{itemize}
\tightlist
\item
  A private key.
\item
  A signature threshold.
\item
  A smart contract condition.
\item
  A role assignment.
\item
  A governance vote.
\item
  A validator majority.
\item
  A sequencer.
\item
  A relayer.
\item
  An oracle committee.
\item
  A backend signing service.
\item
  A custodian's internal approval process.
\item
  A legal administrator with control over off-chain reserves.
\end{itemize}

Authority can be explicit or implicit.

Explicit authority is easy to see:

\begin{lstlisting}
onlyOwner
onlyRole(MINTER_ROLE)
signer == user
threshold signatures >= 3 of 5
\end{lstlisting}

Implicit authority is more dangerous:

\begin{lstlisting}
the frontend constructs the transaction
the indexer decides which position exists
the oracle publishes the price
the backend signs withdrawals
the relayer decides whether to submit messages
the sequencer decides ordering
the deployer can upgrade implementation
the package maintainer can ship malicious code
\end{lstlisting}

Both kinds matter.

\subsection{User Authority}

User authority usually comes from keys, signatures, account-validation
logic, or possession of a capability.

Common user authorities:

\begin{itemize}
\tightlist
\item
  Transfer owned assets.
\item
  Approve a spender.
\item
  Delegate votes.
\item
  Sign a permit.
\item
  Create an order or intent.
\item
  Deposit or withdraw.
\item
  Borrow, repay, liquidate, stake, or claim.
\item
  Configure account recovery or session keys.
\end{itemize}

User authority is often assumed to be benign. That is a mistake. A
malicious user is still a user. A compromised user key is still able to
sign. A tricked user may authorize a harmful transition. A bot may call
public functions in sequences no ordinary user would consider.

The security question is:

\begin{lstlisting}
What can any valid user do, not what should an honest user do?
\end{lstlisting}

\subsection{Admin Authority}

Admin authority is the ability to change protocol behavior, parameters,
or emergency state.

Examples:

\begin{itemize}
\tightlist
\item
  Upgrade implementation.
\item
  Change fees.
\item
  Add collateral types.
\item
  Set risk parameters.
\item
  Pause or unpause.
\item
  Freeze accounts.
\item
  Set oracle addresses.
\item
  Assign minter or burner roles.
\item
  Change bridge validator sets.
\item
  Recover stuck funds.
\item
  Sweep tokens.
\end{itemize}

Admin authority is not automatically bad. Many systems need controlled
administration for upgrades, parameter tuning, emergency response, and
operational recovery. The problem is pretending admin authority is not
trust.

OpenZeppelin's docs make the core point bluntly: access control may
determine who can mint tokens, vote, freeze transfers, or perform other
critical actions.\footnote{OpenZeppelin, ``Access Control'',
  https://docs.openzeppelin.com/contracts/5.x/access-control.} Timelocks
add delay to privileged operations so users can inspect and react before
maintenance actions execute.\footnote{OpenZeppelin, ``Access Control'',
  https://docs.openzeppelin.com/contracts/5.x/access-control.}

Security review must answer:

\begin{itemize}
\tightlist
\item
  What can admins do?
\item
  Can admins directly take user funds?
\item
  Can admins indirectly take user funds by upgrading, changing
  parameters, or replacing dependencies?
\item
  Is there a timelock?
\item
  Who can bypass the timelock?
\item
  Can emergency powers be abused?
\item
  Is there a path to recover if admin keys are lost?
\item
  Are admin actions observable before they take effect?
\end{itemize}

Admin authority should be treated as an attack surface and a trust
assumption, not a footnote.

\subsection{Validator, Miner, Builder, and Sequencer Authority}

Block production authorities do not usually own user assets directly,
but they can affect the history in which asset transitions occur.

Depending on the system, these actors may influence:

\begin{itemize}
\tightlist
\item
  Transaction inclusion.
\item
  Transaction ordering.
\item
  Censorship.
\item
  Reorg risk.
\item
  Finality.
\item
  MEV extraction.
\item
  L2 batch publication.
\item
  Data availability.
\item
  Cross-domain timing.
\end{itemize}

Bitcoin's whitepaper frames the need for a system where participants
agree on a single history of transaction order to address double
spending without a trusted mint.\footnote{Satoshi Nakamoto, ``Bitcoin: A
  Peer-to-Peer Electronic Cash System'',
  https://bitcoin.org/bitcoin.pdf.} That is an authority problem: who
decides which spend came first?

In modern systems, the authority may be distributed across validators,
builders, relays, sequencers, or committees. Even when these actors
cannot forge signatures, they may still decide whether and when valid
transitions become canonical.

Security review must ask:

\begin{lstlisting}
Can block producers reorder this operation profitably?
Can they censor exits or liquidations?
Can they delay oracle updates?
Can they observe intent before inclusion?
Can they create a reorg that invalidates an assumed state?
Can a sequencer include a transaction on L2 before the data or output is secure on L1?
\end{lstlisting}

Authority over ordering is authority over outcomes.

\subsection{Governance Authority}

Governance is privileged execution with a voting interface.

It may control:

\begin{itemize}
\tightlist
\item
  Treasury movement.
\item
  Protocol upgrades.
\item
  Parameter changes.
\item
  Role assignments.
\item
  Emergency powers.
\item
  Validator-set changes.
\item
  Oracle sets.
\item
  Bridge configuration.
\item
  Fee switches.
\end{itemize}

Governance should not be romanticized as ``the community.'' A governance
system is a concrete mechanism with specific state:

\begin{lstlisting}
token balances
delegations
quorum
proposal threshold
voting delay
voting period
timelock delay
execution targets
cancel paths
emergency vetoes
\end{lstlisting}

The real authorities are the actors who can cause proposals to pass and
execute. That may include token whales, delegates, bribery markets,
exchanges, multisigs, guardians, foundations, or attackers using flash
liquidity.

Governance is safer only if its authority is well scoped, observable,
delayed where needed, and resistant to realistic capture.

\subsection{Custodial and Operational Authority}

Some systems depend on people and organizations holding keys or running
processes.

Examples:

\begin{itemize}
\tightlist
\item
  Exchange hot wallets.
\item
  Treasury multisigs.
\item
  MPC signing clusters.
\item
  HSM-backed custodians.
\item
  Backend withdrawal signers.
\item
  Bridge validator operators.
\item
  Oracle publishers.
\item
  Relayers and keepers.
\item
  Deployment engineers.
\item
  CI/CD systems.
\item
  Domain registrars and DNS providers.
\end{itemize}

This authority often lives outside the smart contract but controls
contract-relevant transitions.

For example, a contract may require
\passthrough{\lstinline!onlyWithdrawalSigner!}, but the real question
is:

\begin{lstlisting}
Who can make the signing service sign?
Who can change the signing policy?
Who can deploy signer code?
Who can access production secrets?
Who approves emergency withdrawals?
What happens if one operator is compromised?
What happens if the signer is unavailable?
\end{lstlisting}

Human and organizational authority belongs in the security model.

\subsection{Authority Can Be Delegated, Composed, or Hidden}

Authority is rarely a single direct edge.

Consider:

\begin{lstlisting}
Vault.owner = Timelock
Timelock.proposer = Governor
Governor voting power = token balances + delegation
Token balances include tokens held by an exchange
Timelock.executor = anyone
Guardian can cancel proposals
Proxy admin can upgrade Governor
Frontend helps users delegate
\end{lstlisting}

Who controls the vault?

There is no one-line answer. Control is a graph.

Reviewers should draw that graph:

\begin{lstlisting}
Exchange custody
  -> token voting power
  -> governance outcome
  -> timelock proposal
  -> vault upgrade
  -> user asset movement
\end{lstlisting}

Hidden authority often appears in:

\begin{itemize}
\tightlist
\item
  Upgradeable proxy admins.
\item
  Role admin roles.
\item
  Token allowances.
\item
  Operator approvals.
\item
  Meta-transaction forwarders.
\item
  Signature validators.
\item
  Trusted relayers.
\item
  Oracle owners.
\item
  Mutable registries.
\item
  Off-chain signers.
\item
  Emergency pause modules.
\item
  Dependency contracts that can be replaced.
\end{itemize}

If authority is hidden, the system may appear safer than it is.

\section{Trust Boundaries Across On-Chain, Off-Chain, and Human Systems}

\subsection{A Trust Boundary Marks a Change in Verification}

A trust boundary is where one part of the system must rely on claims,
behavior, or data from another part that it cannot fully verify for
itself.

OWASP describes threat modeling as identifying and understanding threats
and mitigations in the context of protecting something
valuable.\footnote{OWASP, ``Threat Modeling'',
  https://owasp.org/www-community/Threat\_Modeling.} NIST's zero trust
architecture guidance makes a related general security point: trust
should not be implicit merely because of network location or asset
ownership.\footnote{NIST SP 800-207, ``Zero Trust Architecture'',
  https://csrc.nist.gov/pubs/sp/800/207/final.}

For blockchain security, translate that into this rule:

\begin{lstlisting}
Do not trust a component just because it is nearby, official, on-chain,
off-chain, decentralized-looking, or operated by the protocol team.
Name what it can do and what verifies it.
\end{lstlisting}

A trust boundary is crossed when:

\begin{itemize}
\tightlist
\item
  A contract accepts an oracle price.
\item
  A bridge accepts a source-chain proof.
\item
  A rollup accepts a sequencer batch.
\item
  A frontend constructs a transaction for a wallet.
\item
  A backend signs a withdrawal.
\item
  A governance proposal calls an arbitrary target.
\item
  A token contract trusts another token's behavior.
\item
  An indexer supplies state to a UI.
\item
  A user trusts a wallet simulation.
\item
  A protocol trusts a multisig not to abuse an upgrade key.
\end{itemize}

At each boundary, ask:

\begin{lstlisting}
What is being trusted?
Who can lie?
Who verifies?
What happens if the claim is false?
Can the system fail closed?
Can users observe or exit before damage?
\end{lstlisting}

\subsection{On-Chain Does Not Mean Trustless}

On-chain code is transparent and independently executable, but it can
still contain trust assumptions.

On-chain trust assumptions include:

\begin{itemize}
\tightlist
\item
  A privileged role will not abuse power.
\item
  An upgrade admin will not deploy malicious code.
\item
  A governance process will not be captured.
\item
  A token contract will behave like expected.
\item
  A library or dependency address is correct.
\item
  A proxy points to the intended implementation.
\item
  A registry returns correct addresses.
\item
  A timelock delay is long enough for users to react.
\item
  A contract's code is verified and corresponds to the deployed
  bytecode.
\end{itemize}

Even if every line is on-chain, the system may still trust people, keys,
roles, and external state.

On-chain verification reduces some forms of trust. It does not erase
trust.

\subsection{Off-Chain Does Not Mean Insecure}

Off-chain components are not automatically bad. They often provide
necessary services:

\begin{itemize}
\tightlist
\item
  Wallets.
\item
  Frontends.
\item
  Indexers.
\item
  Oracles.
\item
  Keepers.
\item
  Relayers.
\item
  Sequencers.
\item
  Provers.
\item
  Monitoring systems.
\item
  Signing infrastructure.
\item
  Governance discussion forums.
\end{itemize}

The issue is whether their failure modes are explicit and controlled.

For each off-chain component, ask:

\begin{itemize}
\tightlist
\item
  Is it trusted for safety, liveness, or only convenience?
\item
  Can users bypass it?
\item
  Can it forge state transitions, or only suggest them?
\item
  Can it censor or delay?
\item
  Can it lie without being detected?
\item
  Can multiple independent providers reduce risk?
\item
  Is there an on-chain verification step?
\item
  What happens if it disappears?
\end{itemize}

An off-chain indexer that only improves UI performance is one thing. An
off-chain signer that can authorize withdrawals is another. A sequencer
that can delay users but not steal funds is different from a bridge
multisig that can mint wrapped assets.

Do not classify components by location. Classify them by authority and
failure impact.

\subsection{Cross-Chain Systems Create Stacked Trust Boundaries}

Bridges, rollups, appchains, and cross-chain applications combine
multiple state machines. That creates stacked trust boundaries.

A simple lock-and-mint bridge may depend on:

\begin{itemize}
\tightlist
\item
  Source-chain finality.
\item
  Source bridge contract correctness.
\item
  Event or message commitment.
\item
  Relayers.
\item
  Validator or multisig signatures.
\item
  Destination bridge contract correctness.
\item
  Domain separation.
\item
  Replay protection.
\item
  Mint authority for the wrapped asset.
\item
  Governance over bridge configuration.
\end{itemize}

The wrapped token on the destination chain may look like a normal token,
but its backing depends on the entire path.

Security review should avoid statements like:

\begin{lstlisting}
This is backed by locked assets on chain A.
\end{lstlisting}

Instead say:

\begin{lstlisting}
This claim is valid if the source-chain lock event is final,
the bridge verification path is sound,
the relayer or validator assumptions hold,
the destination contract enforces replay protection,
and the bridge authority cannot mint outside the intended rules.
\end{lstlisting}

That is the trust boundary made visible.

\subsection{Human Systems Are Part of the Boundary}

Human systems include:

\begin{itemize}
\tightlist
\item
  Multisig signer coordination.
\item
  Incident response.
\item
  Governance participation.
\item
  Key backup and recovery.
\item
  Deployment review.
\item
  Release management.
\item
  Legal compliance.
\item
  Custody operations.
\item
  Oracle operations.
\item
  Community monitoring.
\end{itemize}

These systems can be strong or weak. They can have redundancy,
separation of duties, checklists, audits, delays, and rehearsed response
plans. Or they can be one engineer with a hot wallet and a production
\passthrough{\lstinline!.env!} file.

Blockchain security often fails because the code was treated as the
whole system.

Ask:

\begin{lstlisting}
Which human process can authorize or block asset movement?
Which process can change code or configuration?
Which process notices malicious changes?
Which process can recover from mistakes?
Which process is assumed to be honest, available, and competent?
\end{lstlisting}

If the answer is ``the team will handle it,'' you have found an
unanalyzed trust boundary.

\subsection{Trust Boundaries Should Be Written as Assumptions}

A good trust-boundary statement is concrete enough to be challenged.

Weak:

\begin{lstlisting}
The oracle is trusted.
\end{lstlisting}

Better:

\begin{lstlisting}
The lending market assumes the oracle reports a price within 1 percent of the executable market price and updates at least every 30 minutes. If the oracle is stale or manipulable, accounts may borrow against inflated collateral and leave bad debt.
\end{lstlisting}

Weak:

\begin{lstlisting}
The multisig is secure.
\end{lstlisting}

Better:

\begin{lstlisting}
The protocol assumes at least 3 of 5 multisig signers remain uncompromised and available. The multisig can upgrade the vault implementation after a 48-hour timelock, which gives users a chance to exit only if withdrawals remain available during the delay.
\end{lstlisting}

Weak:

\begin{lstlisting}
The bridge verifies messages.
\end{lstlisting}

Better:

\begin{lstlisting}
The destination contract accepts messages signed by 13 of 19 bridge validators. If 13 validators collude or are compromised, they can mint unbacked wrapped assets unless a separate rate limit or pause mechanism intervenes.
\end{lstlisting}

This is the style the rest of the book will use.

\section{Mapping Who Can Move, Freeze, Mint, Burn, Upgrade, or Censor Assets}

\subsection{Authority Mapping Is the Practical Output}

An authority map is a structured answer to:

\begin{lstlisting}
Who can do what to which asset, through which transition, under which assumptions?
\end{lstlisting}

This is not a diagramming exercise for its own sake. It is how reviewers
avoid missing the key that can drain the protocol.

A useful authority map includes:

\begin{itemize}
\tightlist
\item
  Asset or state affected.
\item
  Transition or function.
\item
  Required authority.
\item
  Authority holder.
\item
  Delegation path.
\item
  Delay or timelock.
\item
  Revocation path.
\item
  Observability.
\item
  Failure impact.
\item
  Recovery path.
\end{itemize}

\subsection{The Core Authority Verbs}

For each asset, ask who can perform these actions.

{\def\LTcaptype{none} % do not increment counter
\begin{longtable}[]{@{}
  >{\raggedright\arraybackslash}p{(\linewidth - 4\tabcolsep) * \real{0.3333}}
  >{\raggedright\arraybackslash}p{(\linewidth - 4\tabcolsep) * \real{0.3333}}
  >{\raggedright\arraybackslash}p{(\linewidth - 4\tabcolsep) * \real{0.3333}}@{}}
\toprule\noalign{}
\begin{minipage}[b]{\linewidth}\raggedright
Verb
\end{minipage} & \begin{minipage}[b]{\linewidth}\raggedright
Meaning
\end{minipage} & \begin{minipage}[b]{\linewidth}\raggedright
Example risk
\end{minipage} \\
\midrule\noalign{}
\endhead
\bottomrule\noalign{}
\endlastfoot
Move & Transfer or redirect the asset & Unauthorized transfer, malicious
withdrawal \\
Approve & Delegate spending or operation rights & Infinite allowance,
operator abuse \\
Mint & Create new claims & Inflation, unbacked wrapped assets \\
Burn & Destroy claims & User balance destruction, forced redemption \\
Freeze & Prevent movement or use & Censorship, trapped funds \\
Pause & Stop protocol transitions & Emergency control abuse, liveness
failure \\
Upgrade & Change future rules & Malicious implementation, storage
collision \\
Configure & Change parameters or dependencies & Unsafe collateral
factor, bad oracle \\
Liquidate & Seize or sell collateral & Manipulated liquidation,
griefing \\
Settle & Finalize claims or messages & Premature bridge release, wrong
finality \\
Censor & Exclude or delay transitions & Broken exits, liquidation
failure \\
Observe & Learn private or pending intent & MEV, order-flow leakage \\
Recover & Override normal flow after failure & Admin theft disguised as
recovery \\
\end{longtable}
}

This list should feel a little uncomfortable. Good. It reveals how many
ways ``owning'' an asset can be affected by actors other than the
nominal owner.

\subsection{Authority Matrix}

An authority matrix turns the verb list into a review artifact. For a
tokenized vault, a first pass might look like this:

{\def\LTcaptype{none} % do not increment counter
\begin{longtable}[]{@{}
  >{\raggedright\arraybackslash}p{(\linewidth - 14\tabcolsep) * \real{0.1250}}
  >{\raggedright\arraybackslash}p{(\linewidth - 14\tabcolsep) * \real{0.1250}}
  >{\raggedright\arraybackslash}p{(\linewidth - 14\tabcolsep) * \real{0.1250}}
  >{\raggedright\arraybackslash}p{(\linewidth - 14\tabcolsep) * \real{0.1250}}
  >{\raggedright\arraybackslash}p{(\linewidth - 14\tabcolsep) * \real{0.1250}}
  >{\raggedright\arraybackslash}p{(\linewidth - 14\tabcolsep) * \real{0.1250}}
  >{\raggedright\arraybackslash}p{(\linewidth - 14\tabcolsep) * \real{0.1250}}
  >{\raggedright\arraybackslash}p{(\linewidth - 14\tabcolsep) * \real{0.1250}}@{}}
\toprule\noalign{}
\begin{minipage}[b]{\linewidth}\raggedright
Actor or component
\end{minipage} & \begin{minipage}[b]{\linewidth}\raggedright
Move
\end{minipage} & \begin{minipage}[b]{\linewidth}\raggedright
Mint/burn
\end{minipage} & \begin{minipage}[b]{\linewidth}\raggedright
Freeze/pause
\end{minipage} & \begin{minipage}[b]{\linewidth}\raggedright
Upgrade/configure
\end{minipage} & \begin{minipage}[b]{\linewidth}\raggedright
Censor/delay
\end{minipage} & \begin{minipage}[b]{\linewidth}\raggedright
Observe
\end{minipage} & \begin{minipage}[b]{\linewidth}\raggedright
Recover
\end{minipage} \\
\midrule\noalign{}
\endhead
\bottomrule\noalign{}
\endlastfoot
User & Deposit, withdraw, transfer shares & Burn shares on withdrawal &
No & No & No & Own transaction intent & Recover only through wallet/key
process \\
Approved spender & Move approved tokens & No & No & No & No & May infer
user intent from approvals & No \\
Vault contract & Move assets according to rules & Mint/burn shares &
Enforce pause state & No unless upgradeable path exists & Can reject
invalid exits & Emits events and exposes state & May include recovery
hooks \\
Strategist & Move strategy allocation & No & No & Sometimes configure
strategy parameters & Can delay yield movement or withdrawal liquidity &
Sees strategy state & May unwind positions \\
Oracle/reporter & No direct movement & No & No & No & Can delay or
distort valuation-dependent transitions & Observes/report prices & May
trigger fallback modes \\
Guardian & Usually no movement & No & Pause selected transitions &
Sometimes change emergency parameters & Can delay user exits if pause is
broad & Observes incident signals & May trigger recovery mode \\
Governance/timelock & May authorize treasury or strategy movement & May
change mint/burn logic by upgrade & Can assign pausers & Upgrade,
configure, grant roles & Can delay or execute proposals & Observes and
publishes proposals & Can replace roles or implementations \\
Token issuer/custodian & May freeze or redeem underlying token & May
mint/burn underlying token & Can blacklist or freeze & Can change
issuer-side policies & Can block movement of underlying token & Observes
compliance data & May recover or seize under issuer rules \\
\end{longtable}
}

The matrix does not prove the system is safe. It prevents hidden
authority from staying hidden.

\subsection{Example: ERC-20 Token Authority Map}

Consider an ERC-20-like token with minting, burning, pause, and
blacklist features.

{\def\LTcaptype{none} % do not increment counter
\begin{longtable}[]{@{}
  >{\raggedright\arraybackslash}p{(\linewidth - 6\tabcolsep) * \real{0.2500}}
  >{\raggedright\arraybackslash}p{(\linewidth - 6\tabcolsep) * \real{0.2500}}
  >{\raggedright\arraybackslash}p{(\linewidth - 6\tabcolsep) * \real{0.2500}}
  >{\raggedright\arraybackslash}p{(\linewidth - 6\tabcolsep) * \real{0.2500}}@{}}
\toprule\noalign{}
\begin{minipage}[b]{\linewidth}\raggedright
Asset/state
\end{minipage} & \begin{minipage}[b]{\linewidth}\raggedright
Transition
\end{minipage} & \begin{minipage}[b]{\linewidth}\raggedright
Authority
\end{minipage} & \begin{minipage}[b]{\linewidth}\raggedright
Failure mode
\end{minipage} \\
\midrule\noalign{}
\endhead
\bottomrule\noalign{}
\endlastfoot
User balance & \passthrough{\lstinline!transfer!} & Token holder &
Holder key compromise, unsafe recipient assumptions \\
User balance & \passthrough{\lstinline!transferFrom!} & Approved spender
& Infinite allowance, phishing, stale approval \\
Total supply & \passthrough{\lstinline!mint!} & Minter role & Inflation,
unbacked issuance \\
User balance & \passthrough{\lstinline!burnFrom!} & Burner role or
allowance & Unauthorized destruction \\
All transfers & \passthrough{\lstinline!pause!} & Pauser role & Frozen
market, griefing, centralized control \\
Address usability & \passthrough{\lstinline!blacklist!} &
Compliance/admin role & Censorship, mistaken freeze \\
Implementation & \passthrough{\lstinline!upgradeTo!} & Proxy
admin/timelock & Malicious logic, storage corruption \\
Roles &
\passthrough{\lstinline!grantRole!}/\passthrough{\lstinline!revokeRole!}
& Role admin & Privilege escalation \\
\end{longtable}
}

Now ask:

\begin{itemize}
\tightlist
\item
  Who holds each role?
\item
  Can roles be changed?
\item
  Are role changes delayed?
\item
  Are role holders multisigs, EOAs, contracts, or governance?
\item
  Can users detect changes before they are harmed?
\item
  Can the token be safely integrated by protocols that assume standard
  ERC-20 behavior?
\end{itemize}

The ERC-20 standard defines a basic interface, including delegated
spending through approvals and allowances.\footnote{Fabian Vogelsteller
  and Vitalik Buterin, ``ERC-20: Token Standard'',
  https://eips.ethereum.org/EIPS/eip-20.} It does not guarantee that
every token has the same security properties.

\subsection{Example: Vault Authority Map}

Consider a yield vault.

Assets and claims:

\begin{itemize}
\tightlist
\item
  Underlying asset.
\item
  Vault shares.
\item
  Strategy positions.
\item
  Pending withdrawals.
\item
  Reward tokens.
\item
  Admin roles.
\end{itemize}

Authorities:

{\def\LTcaptype{none} % do not increment counter
\begin{longtable}[]{@{}
  >{\raggedright\arraybackslash}p{(\linewidth - 4\tabcolsep) * \real{0.3333}}
  >{\raggedright\arraybackslash}p{(\linewidth - 4\tabcolsep) * \real{0.3333}}
  >{\raggedright\arraybackslash}p{(\linewidth - 4\tabcolsep) * \real{0.3333}}@{}}
\toprule\noalign{}
\begin{minipage}[b]{\linewidth}\raggedright
Action
\end{minipage} & \begin{minipage}[b]{\linewidth}\raggedright
Possible authority
\end{minipage} & \begin{minipage}[b]{\linewidth}\raggedright
Security question
\end{minipage} \\
\midrule\noalign{}
\endhead
\bottomrule\noalign{}
\endlastfoot
Deposit & Any user & Can deposits manipulate share price? \\
Withdraw & Share holder & Can withdrawals be blocked or underpaid? \\
Allocate funds & Strategist or vault logic & Can funds be sent to
malicious strategy? \\
Report profit/loss & Strategy, keeper, or accountant & Can reports
inflate share price? \\
Change strategy & Governance/admin & Is there a timelock and exit
window? \\
Pause deposits & Guardian/admin & Can pause trap users or only reduce
risk? \\
Pause withdrawals & Guardian/admin & Is user exit dependent on trusted
discretion? \\
Upgrade vault & Proxy admin/governance & Can implementation steal
assets? \\
Sweep tokens & Admin/recovery role & Can underlying assets be swept? \\
\end{longtable}
}

The user may ``own'' shares. But the value of those shares depends on
accounting, strategy authority, withdrawal liveness, admin controls,
oracle assumptions, and upgrade authority.

\subsection{Example: Bridge Authority Map}

Consider a bridge that locks tokens on chain A and mints wrapped tokens
on chain B.

{\def\LTcaptype{none} % do not increment counter
\begin{longtable}[]{@{}
  >{\raggedright\arraybackslash}p{(\linewidth - 4\tabcolsep) * \real{0.3333}}
  >{\raggedright\arraybackslash}p{(\linewidth - 4\tabcolsep) * \real{0.3333}}
  >{\raggedright\arraybackslash}p{(\linewidth - 4\tabcolsep) * \real{0.3333}}@{}}
\toprule\noalign{}
\begin{minipage}[b]{\linewidth}\raggedright
Asset/state
\end{minipage} & \begin{minipage}[b]{\linewidth}\raggedright
Authority
\end{minipage} & \begin{minipage}[b]{\linewidth}\raggedright
Boundary
\end{minipage} \\
\midrule\noalign{}
\endhead
\bottomrule\noalign{}
\endlastfoot
Locked source tokens & Source bridge contract & Source-chain execution
and finality \\
Wrapped destination supply & Destination bridge contract & Correct
mint/burn accounting \\
Message validity & Validator set or light client & Signature threshold
or proof verification \\
Message delivery & Relayer network & Liveness and censorship \\
Replay protection & Destination contract state & Nonce/message status
correctness \\
Validator set updates & Governance/admin & Capture or key compromise \\
Emergency pause & Guardian/admin & Abuse vs incident response \\
Upgrade path & Proxy admin/governance & Future rule changes \\
\end{longtable}
}

A wrapped token holder does not merely depend on the token contract.
They depend on every authority in the bridge path.

\subsection{Example: Governance Authority Map}

Governance should be mapped like any other privileged system.

{\def\LTcaptype{none} % do not increment counter
\begin{longtable}[]{@{}
  >{\raggedright\arraybackslash}p{(\linewidth - 4\tabcolsep) * \real{0.3333}}
  >{\raggedright\arraybackslash}p{(\linewidth - 4\tabcolsep) * \real{0.3333}}
  >{\raggedright\arraybackslash}p{(\linewidth - 4\tabcolsep) * \real{0.3333}}@{}}
\toprule\noalign{}
\begin{minipage}[b]{\linewidth}\raggedright
State/action
\end{minipage} & \begin{minipage}[b]{\linewidth}\raggedright
Authority
\end{minipage} & \begin{minipage}[b]{\linewidth}\raggedright
Questions
\end{minipage} \\
\midrule\noalign{}
\endhead
\bottomrule\noalign{}
\endlastfoot
Proposal creation & Token threshold, proposer role & Can attackers
cheaply create proposals? \\
Voting power & Token balances, delegation & Can votes be borrowed,
bribed, or custodied? \\
Quorum & Governance parameters & Can low participation pass major
changes? \\
Proposal outcome & Voting mechanism & Is voting private, delayed,
snapshot-based, or flash-loan-resistant? \\
Execution & Timelock/executor & Who can execute after passing? \\
Cancellation & Guardian, proposer, admin & Can malicious proposals be
stopped? Can legitimate proposals be censored? \\
Target calls & Governor/timelock & Can arbitrary calls seize assets or
grant roles? \\
Governance upgrades & Governance/admin & Can the process change
itself? \\
\end{longtable}
}

Governance is not separate from security. Governance is often the root
authority over the security model.

\subsection{Minimal Authority Mapping Template}

Use this template when starting any review:

\begin{lstlisting}
Asset or state:

Representation:
  Where is it stored or recorded?

Normal holder:
  Who appears to own or control it?

Transitions:
  Which functions, instructions, messages, or processes affect it?

Authorities:
  Who can call or authorize those transitions?

Delegation:
  Can authority be delegated through allowances, operators, roles, signatures, or governance?

Administration:
  Who can change the rules, roles, parameters, or implementation?

Ordering/finality:
  Does the asset depend on inclusion, ordering, confirmations, settlement, or cross-domain finality?

Off-chain dependencies:
  Which oracles, custodians, signers, frontends, relayers, indexers, legal agreements, or human processes matter?

Failure impact:
  What can go wrong if each authority is malicious, compromised, unavailable, or mistaken?

User exit:
  Can users observe danger and leave before damage occurs?
\end{lstlisting}

If the template feels tedious, use it anyway. Most severe bugs hide in
the row someone skipped.

\section{Worked Example: Mapping a Tokenized Vault}

Suppose a protocol offers a vault:

\begin{lstlisting}
Users deposit USDC.
The vault mints shares.
A strategist allocates USDC into yield strategies.
An oracle reports strategy value.
Users can withdraw USDC by burning shares.
Governance can upgrade the vault after a 48-hour timelock.
A guardian can pause deposits and withdrawals immediately.
\end{lstlisting}

The marketing version is:

\begin{lstlisting}
Users own vault shares backed by USDC and strategies.
\end{lstlisting}

The security version begins with assets:

{\def\LTcaptype{none} % do not increment counter
\begin{longtable}[]{@{}ll@{}}
\toprule\noalign{}
Asset & Representation \\
\midrule\noalign{}
\endhead
\bottomrule\noalign{}
\endlastfoot
USDC deposits & Vault token balance \\
Vault shares & Share token balances \\
Strategy assets & External strategy accounting \\
Pending withdrawals & Vault queue state \\
Strategy value & Oracle-reported value \\
Governance power & Governance token/delegation state \\
Emergency authority & Guardian role \\
Upgrade authority & Timelock/governance/proxy admin \\
\end{longtable}
}

Then authority:

{\def\LTcaptype{none} % do not increment counter
\begin{longtable}[]{@{}
  >{\raggedright\arraybackslash}p{(\linewidth - 4\tabcolsep) * \real{0.3333}}
  >{\raggedright\arraybackslash}p{(\linewidth - 4\tabcolsep) * \real{0.3333}}
  >{\raggedright\arraybackslash}p{(\linewidth - 4\tabcolsep) * \real{0.3333}}@{}}
\toprule\noalign{}
\begin{minipage}[b]{\linewidth}\raggedright
Action
\end{minipage} & \begin{minipage}[b]{\linewidth}\raggedright
Authority
\end{minipage} & \begin{minipage}[b]{\linewidth}\raggedright
Risk
\end{minipage} \\
\midrule\noalign{}
\endhead
\bottomrule\noalign{}
\endlastfoot
Deposit & Any user & Share inflation, nonstandard token transfer \\
Withdraw & Share holder & Queue griefing, underpayment, stuck funds \\
Allocate to strategy & Strategist & Loss of funds, malicious strategy \\
Report value & Oracle/reporter & Inflated share price, bad
redemptions \\
Pause deposits & Guardian & Blocks new exposure \\
Pause withdrawals & Guardian & Traps users \\
Upgrade vault & Governance/timelock & Rule replacement \\
Change oracle & Governance/admin & False valuation \\
Change strategist & Governance/admin & Malicious allocation \\
\end{longtable}
}

Now trust boundaries:

\begin{lstlisting}
flowchart LR
    User[User] --> Wallet[Wallet / frontend]
    Wallet --> Vault[Vault contract]
    Vault --> Token[USDC token]
    Vault --> Strategy[Yield strategy]
    Strategy --> Oracle[Oracle / reporter]
    Governance[Governance / timelock] --> Vault
    Guardian[Guardian] --> Vault
    TokenIssuer[USDC issuer] --> Token
\end{lstlisting}

\begin{itemize}
\tightlist
\item
  USDC itself has issuer and blacklist assumptions.
\item
  Strategy value depends on external strategy contracts and
  oracle/reporting.
\item
  Governance depends on voting distribution and timelock execution.
\item
  The guardian is trusted not to freeze withdrawals maliciously.
\item
  Users depend on frontend/wallet correctness to understand approvals.
\item
  The proxy admin path can change the vault's future transition rules.
\end{itemize}

The combined asset/control map is:

{\def\LTcaptype{none} % do not increment counter
\begin{longtable}[]{@{}
  >{\raggedright\arraybackslash}p{(\linewidth - 8\tabcolsep) * \real{0.2000}}
  >{\raggedright\arraybackslash}p{(\linewidth - 8\tabcolsep) * \real{0.2000}}
  >{\raggedright\arraybackslash}p{(\linewidth - 8\tabcolsep) * \real{0.2000}}
  >{\raggedright\arraybackslash}p{(\linewidth - 8\tabcolsep) * \real{0.2000}}
  >{\raggedright\arraybackslash}p{(\linewidth - 8\tabcolsep) * \real{0.2000}}@{}}
\toprule\noalign{}
\begin{minipage}[b]{\linewidth}\raggedright
Asset or control surface
\end{minipage} & \begin{minipage}[b]{\linewidth}\raggedright
Representation
\end{minipage} & \begin{minipage}[b]{\linewidth}\raggedright
Transition
\end{minipage} & \begin{minipage}[b]{\linewidth}\raggedright
Authority
\end{minipage} & \begin{minipage}[b]{\linewidth}\raggedright
Failure impact
\end{minipage} \\
\midrule\noalign{}
\endhead
\bottomrule\noalign{}
\endlastfoot
User deposit & USDC balance held by vault &
\passthrough{\lstinline!deposit!}, \passthrough{\lstinline!withdraw!} &
User, vault logic, USDC token rules & User funds stuck, frozen, or
misaccounted \\
Vault shares & Share token balances & Mint on deposit, burn on
withdrawal & Vault accounting logic & Dilution, unfair redemption, share
inflation \\
Strategy allocation & External strategy position & Allocate, withdraw,
report & Strategist/governance/vault logic & Loss of funds or illiquid
withdrawals \\
Strategy value & Oracle or report state & Report/update valuation &
Oracle/reporter & Overstated assets, unfair share price \\
Withdrawal queue & Vault queue state & Enqueue, process, cancel & User,
vault, keeper/processor & Stuck withdrawals, ordering grief \\
Emergency pause & Pause flag & Pause/unpause & Guardian/governance &
Trapped users or uncontrolled incident \\
Upgrade path & Proxy implementation/admin & Queue/execute upgrade &
Governance/timelock/proxy admin & Complete rule replacement \\
\end{longtable}
}

Now invariants:

\begin{lstlisting}
Shares should represent a proportional claim on vault assets.
Withdrawals should not pay more than the user's fair share.
Deposits should not dilute existing users unfairly.
Strategy accounting should not overstate withdrawable assets.
No privileged role should be able to seize assets without an explicit, documented trust assumption.
Users should be able to exit before delayed malicious upgrades execute, unless the system clearly discloses that they cannot.
\end{lstlisting}

This is the difference between describing a product and reviewing a
system.

\section{Practical Discipline: Replace Nouns With Verbs}

When you see a noun, ask for the verb.

{\def\LTcaptype{none} % do not increment counter
\begin{longtable}[]{@{}
  >{\raggedright\arraybackslash}p{(\linewidth - 2\tabcolsep) * \real{0.5000}}
  >{\raggedright\arraybackslash}p{(\linewidth - 2\tabcolsep) * \real{0.5000}}@{}}
\toprule\noalign{}
\begin{minipage}[b]{\linewidth}\raggedright
Noun
\end{minipage} & \begin{minipage}[b]{\linewidth}\raggedright
Better question
\end{minipage} \\
\midrule\noalign{}
\endhead
\bottomrule\noalign{}
\endlastfoot
Owner & What can this owner do? \\
Admin & Which state transitions can this admin authorize? \\
Guardian & Can the guardian only pause, or can it seize/reconfigure? \\
Governance & Can it call arbitrary targets? Can it upgrade itself? \\
Oracle & Which transitions rely on its data? \\
Relayer & Is it needed for liveness, safety, or both? \\
Sequencer & Can it censor exits or reorder liquidations? \\
Custodian & Which assets or keys does it control? \\
Wallet & What can it cause the user to sign? \\
Bridge & Which domain's state is being trusted by which other domain? \\
\end{longtable}
}

Nouns make systems sound stable. Verbs reveal power.

\section{Review Questions}

\begin{enumerate}
\def\labelenumi{\arabic{enumi}.}
\tightlist
\item
  Why is an on-chain asset better understood as an enforceable claim
  over state than as an object inside an address?
\item
  Give three examples of assets that are not simple token balances.
\item
  What is the difference between legal ownership, protocol authority,
  custody, and economic exposure?
\item
  Why are allowances, operator approvals, upgrade keys, oracle keys, and
  governance roles forms of authority?
\item
  Why can the words ``owner'' and ``account owner'' mislead reviewers
  across Ethereum, Solana, and Move-style systems?
\item
  What is a trust boundary, and why do bridges create stacked trust
  boundaries?
\item
  What information belongs in an authority map?
\item
  Give an example where a user owns an asset economically but cannot
  move it without relying on another actor.
\end{enumerate}

\section{Applied Exercises}

\subsection{Exercise 1: Authority Map a Token}

Choose any token design with:

\begin{itemize}
\tightlist
\item
  Transfers.
\item
  Allowances.
\item
  Minting.
\item
  Burning.
\item
  Pausing.
\item
  Upgradeability.
\end{itemize}

Create an authority map with:

\begin{itemize}
\tightlist
\item
  Asset/state affected.
\item
  Transition.
\item
  Required authority.
\item
  Current authority holder.
\item
  Whether authority is delayed, revocable, or observable.
\item
  Worst-case failure if that authority is compromised.
\end{itemize}

Then answer: which authority is economically equivalent to control over
the whole token?

\subsection{Exercise 2: Trust Boundary and Custody Review}

Analyze both scenarios:

\begin{enumerate}
\def\labelenumi{\arabic{enumi}.}
\tightlist
\item
  A protocol uses an oracle price to decide whether accounts can borrow.
\item
  A user has 100 tokens on a centralized exchange; the exchange displays
  the balance but controls the on-chain wallet keys.
\end{enumerate}

For each scenario, write:

\begin{itemize}
\tightlist
\item
  The asset or claim.
\item
  The actor with protocol authority.
\item
  The custody or reporting dependency.
\item
  The trust-boundary statement.
\item
  The failure impact if the trusted actor is wrong, compromised,
  unavailable, or dishonest.
\end{itemize}

\subsection{Exercise 3: Bridge Claim Analysis}

A wrapped token on chain B claims to represent locked tokens on chain A.
The bridge uses a 7-of-10 validator signature threshold and can be
upgraded by a 3-of-5 multisig.

Answer:

\begin{itemize}
\tightlist
\item
  What is the asset on chain B?
\item
  What is the claimed backing on chain A?
\item
  Who can mint on chain B?
\item
  What trust boundary connects chain A to chain B?
\item
  Which authority is stronger: the validator set or the upgrade
  multisig?
\item
  What happens if 7 bridge validators collude?
\item
  What happens if 3 upgrade multisig signers collude?
\end{itemize}

\subsection{Exercise 4: Replace Nouns With Verbs}

For each phrase, rewrite it as authority over transitions:

\begin{enumerate}
\def\labelenumi{\arabic{enumi}.}
\tightlist
\item
  ``The DAO controls the treasury.''
\item
  ``The guardian protects users.''
\item
  ``The protocol uses a trusted oracle.''
\item
  ``The token is backed by reserves.''
\item
  ``The sequencer runs the rollup.''
\item
  ``The admin can configure markets.''
\end{enumerate}

Example answer style:

\begin{lstlisting}
"The DAO controls the treasury" becomes:
Token holders with enough voting power can pass a proposal that, after the timelock delay, causes the treasury contract to transfer assets or call target contracts.
\end{lstlisting}

\section{Key Takeaways}

\begin{itemize}
\tightlist
\item
  Assets are enforceable claims over state, not physical objects inside
  addresses.
\item
  Authority is the ability to cause accepted state transitions that
  affect assets.
\item
  Ownership, custody, control, legal title, and economic exposure are
  different categories.
\item
  UTXOs, accounts, contract storage, Solana accounts, and Move resources
  encode ownership differently.
\item
  Token balances are incomplete without allowances, operators, roles,
  upgrade paths, and external dependencies.
\item
  Privileged roles, upgrade keys, oracle keys, validator keys, sequencer
  power, and governance votes are assets because they can affect value.
\item
  Trust boundaries appear wherever a system relies on a claim, actor, or
  component it cannot fully verify locally.
\item
  On-chain systems can still contain trust assumptions. Off-chain
  components can be safe if their authority and failure modes are
  constrained.
\item
  Authority maps turn vague security claims into reviewable facts.
\item
  When in doubt, replace nouns with verbs: what can this actor actually
  do?
\end{itemize}

\section{Further Reading}

\begin{itemize}
\tightlist
\item
  ethereum.org,
  \href{https://ethereum.org/developers/docs/accounts/}{``Ethereum
  Accounts''}. Useful for distinguishing externally owned accounts,
  contract accounts, account fields, and wallets.
\item
  Satoshi Nakamoto, \href{https://bitcoin.org/bitcoin.pdf}{``Bitcoin: A
  Peer-to-Peer Electronic Cash System''}. See the transaction and
  double-spend discussion for ownership as signature chain plus agreed
  history.
\item
  Fabian Vogelsteller and Vitalik Buterin,
  \href{https://eips.ethereum.org/EIPS/eip-20}{``ERC-20: Token
  Standard''}. The canonical fungible token interface, including
  transfers, approvals, and allowances.
\item
  William Entriken et al.,
  \href{https://eips.ethereum.org/EIPS/eip-721}{``ERC-721: Non-Fungible
  Token Standard''}. The NFT standard, including token ownership,
  approvals, and operator authority.
\item
  Solana Foundation,
  \href{https://solana.com/docs/core/accounts}{``Accounts''}. The core
  reference for Solana account structure, account ownership, and runtime
  modification rules.
\item
  The Move Book,
  \href{https://move-book.com/reference/abilities/}{``Abilities''}.
  Current reference for Move abilities, linear behavior, storage
  constraints, and resource-like asset modeling.
\item
  OpenZeppelin,
  \href{https://docs.openzeppelin.com/contracts/5.x/access-control}{``Access
  Control''}. Practical reference for ownership, roles, timelocks, and
  privileged smart contract actions.
\item
  OWASP, \href{https://owasp.org/www-community/Threat_Modeling}{``Threat
  Modeling''}. General security framing for identifying threats,
  assumptions, mitigations, and validation.
\item
  NIST SP 800-207,
  \href{https://csrc.nist.gov/pubs/sp/800/207/final}{``Zero Trust
  Architecture''}. Useful general framing for avoiding implicit trust
  and focusing on users, assets, and resources.
\end{itemize}
